% Nouveau EGA V (cote 25): archive cote 25, site page 92 (PDF page 93).
% typescript. Draft compared: none.
% AI-transcribed by Claude Opus 5.5 (claude-opus-5-5, adaptive thinking, effort high), Message Batches API, batch ega25b (msgbatch_01A7n5YuXubUPNK1zsepi8RG), 2026-11-13; checked against both supplied scan parts; see CHECK_DECISIONS.tsv.
\scanpage{25/92}
\mspage{86}

tout au moins sont valables sous d es conditions plus générales que \var{\inserted[typed]{pour}}{pour}{insertion}{ega25_25_092-v0001} la
famille des hyperplans (ou des hypersurfaces de degré donné) de l'espace
projectif. Il me semblerait opportun d'adopter un tel point de vue
axiomatique d\var{\over{an}{è}}{è}{writing over}{ega25_25_092-v0002}s les premiers n$^{\text{o}}$s de la rédaction en forme. Je n'ai pas
tiré au clair pour le moment s'il y a bien moyen de donner une généralisa%
tion dans ce sens de Bertini-Zariski (donc des résultats des n$^{\text{o}}$s 4 et 6),
et j'ai écrit aux compétences (Serre, Zariski) pour leur demander s'ils
avaient connaissance d'une telle extension; j'ai bien l'impression en
effet que des hypothèses de nature différentielle simple, du genre de cel
% NOTE: the end of the line runs off the right edge of the scan after "cel"
les envisagées ci-dessous, doivent suffire à impliquer Bertini-Zariski.
Si les compétences ne peuvent nous informer de façon satisfaisante, il
faudrait essayer de tirer la chose au clair par nos propres moyens.

On part d'un diagramme \var{\inserted[typed]{commutatif}}{commutatif}{insertion}{ega25_25_092-v0003} de morphismes de présentation finie

\var{\margin{(D)}}{\placedmargin{(D)}}{marginal note placed}{ega25_25_092-v0004}
\diagram{\[
\begin{array}{ccc}
P & \hand{\longleftarrow} & \underline{P} \\
\hand{\vert} & & \hand{\downarrow} \\
S & \hand{\longleftarrow} & G
\end{array}
\]}

(dans le cas d'application principal, P est un fibré projectif, G une
grasm\unsure{a|à}nienne dudit, $\underline{P}$ le préschéma d'incidence). Dans lescas lesplus
importants, le morphisme correspondant $\underline{P} \hand{\longrightarrow} Px_S G$ sera une immersion
fermée, et on considère G comme \var{\xout{définissant une}}{}{typewriter strike-out}{ega25_25_092-v0005} \var{\inserted[typed]{schéma des paramètres d'une}}{schéma des paramètres d'une}{insertion}{ega25_25_092-v0006} familled e sous-
préschémas fermés des fibres de P sur S, d e façon précise si \hand{\unsure{\xi|\zeta}} \hand{$\in$} G,
$\underline{P}_{\hand{\unsure{\xi|\zeta}}}$ est un sous-préschéma fermé de $P_{s\ k(\hand{\unsure{\xi|\zeta}})}$ , où s est le point de S en
dessous de \hand{\unsure{\xi|\zeta}} .(D'ailleurs, pour la plupart des énoncés dans ce contexte,
on aura sans doute S=Spec(k)\hand{)}. Dans le cas général, on peut encore con-%
sidérer G comme schéma des paramètres d'une famille de préschémas
au dessus des fibres de P sur S, à \hand{\unsure{\xi|\zeta}} correspondant $\underline{P}_{\hand{\unsure{\xi|\zeta}}}$ au dessus de
$P_{s\ k(\hand{\unsure{\xi|\zeta}})}$ . Bien entendu, au lieu de prendre pour \hand{\unsure{\xi|\zeta}} un point \var{\xout{de}}{}{typewriter strike-out}{ega25_25_092-v0007} (absolu)

\var{\margin{92}}{\placedmargin{92}}{marginal note placed}{ega25_25_092-v0008}
